\chapter{Resultados}\label{cap_4}

Este capítulo apresenta os resultados em uma ordem que acompanha o raciocínio da investigação, e não a cronologia dos experimentos. Primeiro a reprodução do trabalho de referência, que define o ponto de partida. Depois a avaliação dos critérios de \ac{AL} sobre a seleção de imagens, que é a forma usual de aplicar a técnica. Em seguida o diagnóstico que mede quanta margem existe na escolha da região, e a caracterização de para onde cada critério aponta, que juntos explicam por que os critérios conhecidos não alcançam essa margem. Os componentes de \ac{CL} e \ac{IL} vêm na sequência. O capítulo fecha com a intervenção derivada do diagnóstico e com os resultados de natureza metodológica.

Todos os números vêm do registro de execuções descrito no Capítulo \ref{cap_3}. Cada tabela é gerada a partir desse registro por um programa, sem passo manual, e acompanha um arquivo de procedência que lista os identificadores das execuções que produziram cada linha. Nenhum valor foi transcrito à mão.

\section{Reprodução do trabalho de referência}

A reprodução do arranjo descrito por \cite{soares2026flim} confirma o comportamento central relatado pelos autores: o Fβ cresce de forma monotônica à medida que imagens são acrescentadas ao conjunto de treino e satura entre três e quatro imagens, atingindo valores próximos aos publicados para o codificador.

Duas diferenças de implementação precisam ser declaradas antes de qualquer comparação numérica. A primeira é que a avaliação usa limiarização de Otsu seguida de filtro de área, e não o delineamento por árvores dinâmicas empregado no trabalho original, cujo executável é um binário compilado para Linux que não roda no ambiente usado aqui. A segunda é que a seleção de imagens do Algoritmo 1 do trabalho de referência exige o ground truth de todo o conjunto disponível para escolher as poucas imagens que serão anotadas, o que os próprios autores declaram ao descrever a abordagem como supervisionada.

Essa segunda diferença não é um detalhe de implementação, e sim o ponto de partida desta dissertação. Se a escolha das imagens depende de conhecer a resposta em todas elas, o procedimento não é aplicável na situação que motiva o \ac{FLIM}, que é justamente a ausência de anotação. A pergunta que organiza os experimentos seguintes é quanto do benefício dessa escolha se recupera com critérios que não usam ground truth algum.

A reimplementação do Algoritmo 1 revelou ainda duas propriedades não discutidas no trabalho original. O critério de selecionar a imagem de pior Fβ tende a escolher imagens insegmentáveis, e não imagens informativas: as imagens com Fβ igual a zero são aquelas cujas componentes previstas caem fora da faixa aceita pelo filtro de área, de modo que o algoritmo prioriza casos em que nenhuma anotação ajudaria. Além disso, sem uma lista de imagens rejeitadas, a imagem removida pelo retrocesso volta a ser a de menor Fβ na iteração seguinte, e o procedimento pode entrar em ciclo.

\section{Seleção de imagens}

A primeira pergunta é se algum critério de \ac{AL} sem ground truth supera a escolha aleatória das imagens a anotar. Foram avaliados critérios das duas famílias predominantes na literatura: incerteza, representada por entropia, \textit{least confidence}, margem e BALD; e diversidade, representada por CoreSet, BADGE e medoide. A referência superior é o oracle, que usa o ground truth para escolher a imagem de pior desempenho, e a referência inferior é o sorteio.

A Tabela \ref{tab_geral_al} resume a comparação nos três conjuntos de dados, para quatro orçamentos de anotação.

\begin{table}[htb]
\centering
\caption{Critérios de Active Learning por orçamento e conjunto de dados.}
\label{tab_geral_al}
% Active Learning no FLIM: todos os critérios, todos os K (FLIM_lm)
% GERADO por flim_al/tabelas.py em 2026-10-06T16:49:56+00:00
% NAO EDITE A MAO: a proxima geracao sobrescreve, e o numero editado deixa de bater com evidencia/execucoes/.
% Proveniencia: geral_al_lm.proveniencia.json
\begin{tabular}{llrrrrrr}
\hline
critério & K & Parasitas & BraTS & Conjuntivite & $\Delta$ médio vs sorteio & p & vence o sorteio em \\
\hline
CoreSet (imagem) & 2 & 0.543 & 0.000 & 0.642 & -0.061 & 0.5568 & 3 de 9 \\
CoreSet (imagem) & 3 & 0.538 & 0.033 & 0.587 & 0.023 & 0.6074 & 3 de 9 \\
CoreSet (imagem) & 5 & 0.640 & 0.489 & 0.748 & 0.112 & 0.3323 & 7 de 9 \\
CoreSet (imagem) & 8 & 0.690 & 0.230 & 0.636 & -0.067 & 0.4144 & 3 de 9 \\
Entropia (imagem) & 2 & 0.457 & 0.001 & 0.635 & -0.092 & 0.3677 & 3 de 9 \\
Entropia (imagem) & 3 & 0.474 & 0.000 & 0.591 & -0.008 & 0.9104 & 2 de 9 \\
Entropia (imagem) & 5 & 0.599 & 0.498 & 0.667 & 0.074 & 0.5274 & 6 de 9 \\
Entropia (imagem) & 8 & 0.681 & 0.000 & 0.604 & -0.157 & 0.0427 & 2 de 9 \\
Least confidence (imagem) & 2 & 0.420 & 0.001 & 0.635 & -0.105 & 0.3115 & 3 de 9 \\
Least confidence (imagem) & 3 & 0.451 & 0.000 & 0.591 & -0.016 & 0.8178 & 2 de 9 \\
Least confidence (imagem) & 5 & 0.583 & 0.498 & 0.667 & 0.069 & 0.5531 & 6 de 9 \\
Least confidence (imagem) & 8 & 0.684 & 0.000 & 0.441 & -0.210 & 0.0327 & 2 de 9 \\
AL de REGIAO & 2 & 0.142 & 0.252 & 0.442 & -0.178 & 0.2461 & 1 de 9 \\
AL de REGIAO & 3 & 0.280 & 0.257 & 0.443 & -0.036 & 0.7332 & 3 de 9 \\
AL de REGIAO & 5 & 0.039 & 0.499 & 0.261 & -0.247 & 0.1570 & 2 de 9 \\
AL de REGIAO & 8 & 0.412 & 0.525 & 0.308 & -0.170 & 0.1903 & 2 de 9 \\
Oracle (usa GT) & 2 & 0.093 & 0.196 & 0.274 & -0.269 & 0.0921 & 1 de 9 \\
Oracle (usa GT) & 3 & 0.420 & 0.630 & 0.312 & 0.091 & 0.5347 & 5 de 9 \\
Oracle (usa GT) & 5 & 0.618 & 0.445 & 0.572 & 0.031 & 0.7792 & 6 de 9 \\
Oracle (usa GT) & 8 & 0.518 & 0.008 & 0.535 & -0.232 & 0.0107 & 0 de 9 \\
\hline
\end{tabular}

\end{table}

% Nota de geral_al_lm. GERADA por flim_al/tabelas.py.
\footnotesize Uma linha por (critério, K). As colunas de dataset são o F$\beta$ médio sobre as sementes. \textbf{\texttt{FLIM puro (sorteio)} é a referência}: todos os braços recebem o MESMO número de imagens e são medidos no mesmo conjunto de teste. Os critérios de IMAGEM escolhem quais imagens anotar; \textbf{\texttt{AL de REGIÃO} usa as mesmas imagens do sorteio e muda só ONDE o traço cai}, com a mesma contagem de pixels --- o $\Delta$ dele mede posição de anotação, não escolha de imagem. \textbf{\texttt{Oracle} usa o gabarito} para escolher a imagem em que o modelo vai pior (passo 7 do Algoritmo 1 do artigo): não é método, é referência, e o fato de ele também não se destacar é o achado mais informativo daqui. $\Delta$ e p são pareados por semente dentro de cada dataset e K, depois agregados. \texttt{vence o sorteio em} conta as células (dataset $\times$ semente) favoráveis. Decoder FLIM\_lm.
\normalsize


Nenhum critério apresenta superioridade sobre o sorteio que sobreviva à correção de Benjamini-Hochberg. O resultado se mantém nos três conjuntos de dados e nos quatro orçamentos. Vale notar que o oracle, que lê o ground truth e portanto representa o teto do que a seleção de imagens poderia entregar, também não se destaca de forma consistente, perdendo do sorteio em parte considerável das células. Isso sugere que, no regime de três a cinco imagens, a identidade das imagens escolhidas importa menos do que a formulação do Algoritmo 1 pressupõe.

É importante enunciar esse resultado com precisão. O que foi observado é ausência de evidência de superioridade nas condições avaliadas, e não demonstração de equivalência. Os intervalos de confiança são compatíveis com efeitos pequenos em qualquer direção, e o número de repetições limita o tamanho de efeito detectável.

\section{Seleção de região: a margem existe}

A ausência de ganho na seleção de imagens admite duas leituras opostas. Pode ser que a escolha não importe, caso em que a técnica seria inaplicável por natureza. Ou pode ser que a escolha importe e os critérios não a capturem, caso em que o problema permanece aberto. Distinguir as duas exige medir o teto, e não apenas o desempenho.

O experimento de ganho marginal fixa a partição, as imagens de treino, os marcadores iniciais, o codificador de partida e o conjunto de validação, e varia uma única coisa: qual região recebe os aproximadamente trezentos pixels seguintes de anotação. Para cada região candidata o codificador é reestimado e o Fβ é medido novamente, de modo que a diferença entre candidatos isola o efeito da posição da anotação.

\begin{table}[htb]
\centering
\caption{Ganho marginal de uma anotação adicional, por tipo de região escolhida.}
\label{tab_ganho_marginal}
% Ganho marginal de um marcador adicional — Parasitas, FLIM_lm, base funcional
% GERADO por flim_al/tabelas.py em 2026-10-06T16:49:57+00:00
% NAO EDITE A MAO: a proxima geracao sobrescreve, e o numero editado deixa de bater com evidencia/execucoes/.
% Proveniencia: ganho_marginal_schisto_FLIM_lm_funcional.proveniencia.json
\begin{tabular}{lrrrrrrr}
\hline
o que recebeu o marcador & n (sementes) & $\Delta$ F$\beta$ médio & IC95\% & pior & melhor & $\Delta$ vs sorteado & p \\
\hline
melhor região (teto amostrado) & 10 & 0.2191 & [0.0851, 0.3530] & 0.0300 & 0.5777 & 0.3089 & 0.0000 \\
argmax entropia (o AL) & 10 & 0.0843 & [-0.0577, 0.2264] & -0.1196 & 0.4676 & 0.1741 & 0.0157 \\
argmax least confidence & 10 & 0.1078 & [-0.0267, 0.2423] & -0.1168 & 0.4647 & 0.1976 & 0.0071 \\
candidato sorteado (média) & 10 & -0.0898 & [-0.1697, -0.0099] & -0.2672 & 0.1461 & --- & --- \\
pior região & 10 & -0.3490 & [-0.4764, -0.2216] & -0.6091 & -0.0433 & -0.2592 & 0.0000 \\
imagem nova inteira (orçamento MAIOR) & 10 & 0.0506 & [-0.0402, 0.1414] & -0.1427 & 0.3172 & 0.1404 & 0.0002 \\
sorteio uniforme (reponderado, sem GT) & 10 & -0.1050 & [-0.1789, -0.0311] & -0.3038 & 0.1034 & -0.0152 & 0.1930 \\
\hline
\end{tabular}

\end{table}

% Nota de ganho_marginal_schisto_FLIM_lm_funcional. GERADA por flim_al/tabelas.py.
\footnotesize $\Delta$ = F$\beta$ na \textbf{validação} com o marcador extra menos F$\beta$ sem ele. Dentro de cada semente, a partição, as imagens iniciais, os marcadores base, o encoder inicial e o conjunto de validação são idênticos --- a única variável é qual região recebeu os ~300 px adicionais. O treino é determinístico \textbf{dentro de um processo}, a partir dos mesmos arquivos de marcador (verificado: três repetições dão F$\beta$ idêntico até a décima casa). \textbf{Entre} execuções há ruído: a reexecução completa do schisto deu F$\beta$ idêntico em 538 de 540 registros, e nos 2 restantes o desvio máximo foi 0,0034 --- duas ordens de grandeza abaixo dos efeitos reportados. \textbf{n é o número de sementes}: os candidatos de uma mesma semente compartilham encoder e são correlacionados. \texttt{melhor região (teto)} escolhe pelo F$\beta$ da validação e \textbf{não é uma estratégia} --- mede o que existe para capturar, \textbf{entre os 22 candidatos examinados} de ~380 superpixels; não é teto absoluto. Linhas restritas a: \textbf{base funcional}. A separação por regime é post-hoc e forçada pela aritmética --- $\Delta$ a partir de F$\beta$=0 exato não pode ser negativo, então as duas populações não têm média comparável. As sementes compartilham o pool e o conjunto de validação: o IC95\% vale para \textbf{esta} validação sob sorteio das imagens de treino, e não generaliza para o dataset. O conjunto de teste não foi tocado por esta campanha.
\normalsize


A dispersão entre candidatos é grande. A melhor região examinada supera a região sorteada por margens entre 0,133 e 0,303 de Fβ, com significância nas seis células formadas por conjunto de dados e decodificador. A pior região, por sua vez, fica entre 0,199 e 0,502 abaixo do sorteio, também nas seis células. Uma única anotação adicional, mal posicionada, chega a levar um codificador de Fβ 0,78 para zero.

Dois esclarecimentos são necessários. A linha identificada como melhor região escolhe o candidato pelo Fβ de validação e portanto não constitui uma estratégia aplicável: ela mede o que existe para ser capturado, e não um método. E o valor medido é um máximo sobre os candidatos examinados, não um teto absoluto, já que apenas parte dos superpixels de cada imagem foi avaliada.

Ainda assim, o resultado responde à ambiguidade. Existe margem considerável na escolha da região, e ela é da ordem de grandeza que separa um modelo mediano de um modelo bom. A pergunta passa a ser por que os critérios não a alcançam.

\section{Para onde os critérios apontam}

A resposta está em medir o comportamento dos critérios, e não apenas o seu desempenho. Como escolher uma região exige apenas pontuar e ordenar, essa medição dispensa reestimar o codificador para cada candidato e pode ser feita sobre dez repetições nos três conjuntos de dados a custo baixo.

A fração de objeto de uma região descreve onde ela está: valor próximo de zero indica fundo puro, valor próximo de um indica interior do objeto, e valores intermediários indicam que a região atravessa a fronteira. Essa distinção importa porque a comparação com os marcadores reais dos especialistas, apresentada na Tabela \ref{tab_artigo_regiao}, mede que a anotação de borda vale entre 0,126 e 0,166 de Fβ sobre a anotação em qualquer outro ponto interno ao objeto.

\begin{table}[htb]
\centering
\caption{Anotação dos especialistas comparada a cliques realocados.}
\label{tab_artigo_regiao}
% Tabela do artigo × cliques realocados por AL × realocados ao acaso
% GERADO por flim_al/tabelas.py em 2026-10-08T14:01:04+00:00
% NAO EDITE A MAO: a proxima geracao sobrescreve, e o numero editado deixa de bater com evidencia/execucoes/.
% Proveniencia: artigo_vs_regiao.proveniencia.json
\begin{adjustbox}{max width=\linewidth}
\begin{tabular}{lrrrrrrrrrrr}
\toprule
decoder & n & F$\beta$ artigo & F$\beta$ AL região & F$\beta$ aleatório & $\Delta$ AL-artigo & p & $\Delta$ aleat-artigo & p  & $\Delta$ AL-aleat & p   & F$\beta$ publicado (A) \\
\midrule
FLIM\_lm & 6 & 0.744 & 0.622 & 0.669 & -0.122 & 0.0370 & -0.075 & 0.0480 & -0.047 & 0.3088 & 0.860 \\
FLIM\_pb & 6 & 0.720 & 0.494 & 0.537 & -0.225 & 0.0140 & -0.182 & 0.0240 & -0.043 & 0.4223 & 0.857 \\
FLIM\_mb & 6 & 0.718 & 0.502 & 0.544 & -0.217 & 0.0106 & -0.174 & 0.0121 & -0.042 & 0.3997 & 0.843 \\
FLIM\_at & 6 & 0.518 & 0.303 & 0.354 & -0.214 & 0.0051 & -0.164 & 0.0015 & -0.050 & 0.1178 & 0.740 \\
FLIM\_ts & 6 & 0.591 & 0.466 & 0.491 & -0.124 & 0.0031 & -0.099 & 0.0000 & -0.025 & 0.2922 & 0.747 \\
FLIM\_lt & 6 & 0.649 & 0.549 & 0.575 & -0.100 & 0.0043 & -0.075 & 0.0238 & -0.026 & 0.2467 & 0.810 \\
FLIM\_ts* & 6 & 0.506 & 0.349 & 0.391 & -0.158 & 0.0029 & -0.115 & 0.0005 & -0.042 & 0.1427 & --- \\
\bottomrule
\end{tabular}
\end{adjustbox}

\end{table}

% Nota de artigo_vs_regiao. GERADA por flim_al/tabelas.py.
\footnotesize Mesmas imagens --- as que os especialistas A e B anotaram ---, mesmo número de cliques e mesmo traço de fundo (copiado, não regerado) nos três braços. Só muda para onde vão os cliques de objeto. \textbf{Não há seleção de imagem: é seleção de região.} As três diferenças decompõem o efeito --- \texttt{aleat-artigo} isola o ato de tirar o clique da borda, e \texttt{AL-aleat} isola a escolha do Active Learning. Sem a terceira, um $\Delta$ negativo na primeira seria creditado ao AL quando pode ser todo da segunda. Cada célula é um par (usuário, split), n = 6, pareado. \texttt{F$\beta$ publicado (A)} vem da Tabela III do artigo e \textbf{não} é comparável com a coluna reproduzida: o artigo aplica Dynamic Trees, cujo binário não executa nesta máquina. Avaliação em 250 imagens de Z$_1$\textbackslash{}T, as mesmas em todos os braços.
\normalsize


A medição do comportamento dos critérios produz três constatações.

A primeira é que os critérios de incerteza encontram o objeto. Regiões de interior puro representam de 2\% a 8\% do conjunto de candidatos, e entropia e \textit{least confidence} selecionam interior em 40\% a 100\% das escolhas, o que é um enriquecimento expressivo sobre a taxa de base.

A segunda é que eles erram a geometria. As escolhas atravessam a fronteira em 0\% a 20\% dos casos, embora seja justamente a travessia que a comparação com os especialistas mostra ter valor.

A terceira é a mais consequente: regiões que atravessam a fronteira representam apenas 2\% a 12\% do conjunto de candidatos. A anotação valiosa quase não está disponível para ser escolhida. Isso não é uma falha dos critérios, e sim do procedimento que gera a lista. O SLIC é construído para que as fronteiras dos superpixels coincidam com as bordas presentes na imagem, de modo que seus segmentos tendem a ficar inteiramente dentro ou inteiramente fora do objeto. A propriedade que faz dele um bom algoritmo de superpixels é a mesma que o torna inadequado como gerador de candidatos para anotação de borda.

Os critérios de diversidade, avaliados no nível de região pela primeira vez neste trabalho, comportam-se de forma ainda menos favorável. CoreSet localiza o objeto apenas no conjunto de conjuntivite, onde os objetos são grandes, e nos demais seleciona fundo em todas as escolhas, sem qualquer enriquecimento sobre a taxa de base. O medoide nunca localiza o objeto. A explicação é o desequilíbrio do conjunto de candidatos: a diversidade é calculada sobre a distribuição de aparência das regiões, e quando o objeto ocupa cerca de 3\% da imagem ele não participa dessa distribuição de forma relevante. Nos dois conjuntos de objetos pequenos, CoreSet e medoide, que são critérios opostos, produzem a mesma resposta.

\section{Aprendizado contínuo sem retropropagação}

O esquecimento catastrófico, tal como descrito na Seção \ref{sec_cl}, é atribuído ao deslocamento dos pesos durante o ajuste por gradiente, e os métodos que o combatem atuam sobre esse deslocamento. Como o \ac{FLIM} não possui otimizador, nenhum desses métodos se aplica diretamente. O fenômeno, contudo, existe, e os resultados da seção anterior o quantificam: uma anotação adicional pode reduzir o Fβ em até meio ponto.

A inspeção do procedimento de estimação dos filtros localiza a origem. Os candidatos a filtro extraídos de todas as imagens e de todos os marcadores são reunidos em um único conjunto, que então é reduzido ao número de canais da camada por um agrupamento. Acrescentar marcadores acrescenta candidatos, o agrupamento é refeito sobre o conjunto ampliado, e os centróides que estavam em uso se deslocam. O banco de filtros é compartilhado por todas as imagens, inclusive as não anotadas, de modo que a perturbação é global e não local.

O mecanismo de retenção proposto no Capítulo \ref{cap_3} interpola o banco: cada centróide novo é ancorado no centróide anterior mais próximo, com um parâmetro de plasticidade que vai de retenção total a reestimação completa. A verificação de corretude confirma as duas extremidades, reproduzindo o comportamento original em um extremo e impedindo a entrada de qualquer direção nova no outro.

A avaliação, porém, não encontrou valor de plasticidade que convertesse o mecanismo em ganho mensurável. O teste primário, declarado antes da execução, comparou plasticidade intermediária contra reestimação completa e não atingiu significância, com sinais inconsistentes entre os tipos de região examinados. A retenção total chegou a parecer vantajosa em uma análise exploratória, mas o resultado não sobreviveu ao aumento do número de repetições nem ao teste em imagens inéditas, e o critério de falseamento declarado no pré-registro foi atingido: a qualidade final sob retenção total é menor, o que contraria a explicação proposta de que congelar o banco evitaria dano sem custo.

\section{Esforço de anotação}

O eixo da segmentação interativa é diferente do Fβ a orçamento fixo. A métrica usual é o número de cliques necessários para que a qualidade ultrapasse um alvo, com um teto de interações. Um método pode não alterar o Fβ a orçamento fixo e ainda assim reduzir o esforço de interação.

\begin{table}[htb]
\centering
\caption{Cliques até a qualidade alvo, por plasticidade do banco de filtros.}
\label{tab_noc}
% Cliques até a IoU alvo, por plasticidade do banco de filtros
% GERADO por flim_al/tabelas.py em 2026-10-08T14:01:07+00:00
% NAO EDITE A MAO: a proxima geracao sobrescreve, e o numero editado deixa de bater com evidencia/execucoes/.
% Proveniencia: noc_noc.proveniencia.json
\begin{adjustbox}{max width=\linewidth}
\begin{tabular}{lrrrrrrrrr}
\toprule
dataset & $\alpha$ & n & NoC médio & NoC mediano & taxa de sucesso & IoU final & $\Delta$ NoC vs $\alpha$=1 & p & McNemar \\
\midrule
Parasitas & 0.00 & 38 & 11.03 & 12.0 & 13\% & 0.5267 & -0.53 & 0.1031 & 3/0 (p=0.250) \\
Parasitas & 0.50 & 22 & 11.09 & 12.0 & 14\% & 0.5490 & -0.55 & 0.2876 & 2/0 (p=0.500) \\
Parasitas & 1.00 & 38 & 11.55 & 12.0 & 5\% & 0.5501 & --- & --- & --- \\
BraTS & 0.00 & 40 & 7.42 & 12.0 & 48\% & 0.4096 & -0.35 & 0.5829 & 4/3 (p=1.000) \\
BraTS & 0.50 & 20 & 6.00 & 3.0 & 60\% & 0.5273 & 0.05 & 0.7715 & 1/1 (p=1.000) \\
BraTS & 1.00 & 40 & 7.78 & 12.0 & 45\% & 0.4064 & --- & --- & --- \\
Conjuntivite & 0.00 & 1 & 12.00 & 12.0 & 0\% & 0.0000 & --- & --- & 0/1 (p=1.000) \\
Conjuntivite & 0.50 & 1 & 6.00 & 6.0 & 100\% & 0.7617 & --- & --- & 0/0 (p=1.000) \\
Conjuntivite & 1.00 & 1 & 6.00 & 6.0 & 100\% & 0.7764 & --- & --- & --- \\
\bottomrule
\end{tabular}
\end{adjustbox}

\end{table}

% Nota de noc_noc. GERADA por flim_al/tabelas.py.
\footnotesize NoC@X é o número de cliques até a IoU passar do alvo, com teto de 12; é o eixo de RITM (arXiv:2102.06583) e SimpleClick (arXiv:2210.11006), e \textbf{não} se compara com o NoC publicado deles, que usa outros conjuntos e outra definição de alvo. Imagem que não atinge o alvo entra com o teto, e por isso a \textbf{taxa de sucesso} aparece em toda linha: um braço pode mostrar NoC baixo por desistir mais cedo. \texttt{$\alpha$=1} é o FLIM puro, que refaz o banco de filtros a cada clique; \texttt{$\alpha$=0} congela o banco, e o clique chega ao modelo só pelos rótulos que o decoder \texttt{labeled\_marker} lê. Os testes são \textbf{separados por dataset}: agregar Schisto e BraTS, com 7\% e 30\% de sucesso, inflou um achado exploratório a p=0,0265 que separado não passa de p=0,10. O clique é simulado do ground truth pelo protocolo de Xu et al. (2016), com usuário que acerta sempre, então a tabela mede \textbf{número de interações e não tempo de especialista}.
\normalsize


O comportamento mais informativo dessa avaliação não é o valor médio, e sim a ocorrência de regressão. Em várias execuções o modelo piora ao longo da sequência de cliques: em um caso representativo, doze interações levaram a qualidade de 0,724 para 0,702. O especialista anota repetidamente e o modelo não avança, porque cada anotação nova reestima o banco de filtros e desfaz parte do que havia sido construído.

A proteção do banco não reduziu o número de cliques de forma significativa, e o confirmatório em imagens inéditas falseou tanto a hipótese quanto o mecanismo proposto para ela. A taxa de sucesso é reportada ao lado da contagem de cliques porque imagens que não atingem o alvo entram com o teto, e um braço poderia exibir contagem baixa simplesmente por desistir antes.

\section{A intervenção no gerador de candidatos}

As seções anteriores convergem para um diagnóstico. Existe margem considerável na escolha da região. Os critérios de incerteza localizam o objeto mas pousam no interior. Os critérios de diversidade não localizam o objeto. E a anotação que os marcadores dos especialistas mostram ser valiosa, a que atravessa a fronteira, corresponde a uma fração mínima do conjunto de candidatos. Como um critério apenas ordena a lista que recebe, nenhuma escolha de escore poderia resolver a limitação.

A intervenção que decorre desse diagnóstico atua sobre o gerador. Em lugar de segmentar a imagem em superpixels, os candidatos passam a ser discos centrados em pontos do contorno previsto pelo modelo, de modo que atravessam a fronteira por construção. O contorno usado é o previsto, não o do ground truth, e portanto o procedimento é aplicável.

O experimento fixa o critério de seleção em sorteio nos dois braços e mantém o orçamento de anotação idêntico, verificado em mediana de trezentos pixels em ambos. A única diferença entre os braços é a origem da lista de candidatos, o que isola o efeito do gerador.

\begin{table}[htb]
\centering
\caption{Troca do gerador de candidatos, com critério e orçamento fixos.}
\label{tab_gerador}
% Gerador de candidatos: contorno previsto contra superpixel
% GERADO por flim_al/tabelas.py em 2026-10-08T14:01:07+00:00
% NAO EDITE A MAO: a proxima geracao sobrescreve, e o numero editado deixa de bater com evidencia/execucoes/.
% Proveniencia: gerador_contorno.proveniencia.json
\begin{adjustbox}{max width=\linewidth}
\begin{tabular}{lrrrrrrrrrrr}
\toprule
dataset & n & $\Delta$ médio & IC95\% & vitórias & p (t) & p (sinais) & p (perm.) & $\Delta$ pior caso & p  & $\Delta$ dispersão & p   \\
\midrule
Parasitas & 10 & 0.0772 & [0.0393, 0.1151] & 10 de 10 & 0.0013 & 0.0020 & 0.0020 & 0.2153 & 0.0001 & -0.1091 & 0.0000 \\
BraTS & 3 & -0.0089 & [-0.3263, 0.3084] & 1 de 3 & 0.9145 & 1.0000 & 0.7500 & -0.1063 & 0.2703 & 0.0464 & 0.1376 \\
Conjuntivite & 5 & 0.0452 & [-0.0452, 0.1357] & 4 de 5 & 0.2374 & 0.3750 & 0.2500 & 0.1759 & 0.0480 & -0.0577 & 0.1492 \\
\bottomrule
\end{tabular}
\end{adjustbox}

\end{table}

% Nota de gerador_contorno. GERADA por flim_al/tabelas.py.
\footnotesize $\Delta$ é a variação de F$\beta$ na validação ao acrescentar \textbf{uma} anotação, do braço de contorno menos o braço de superpixel. Os dois sorteiam o candidato e gastam os mesmos ~300 px: a \textbf{única} diferença é o gerador da lista, e por isso a comparação isola o gerador e não diz nada sobre qual escore usar. A unidade é a \textbf{semente}; os dois decodificadores de uma semente compartilham o encoder e entram colapsados por média. Pior caso e dispersão foram declarados no pré-registro, com o mecanismo escrito antes de rodar: o candidato de contorno é pequeno e local, acrescenta menos filtros, e o ganho marginal mede correlação negativa entre filtros acrescentados e ganho. O nulo do BraTS é \textbf{previsto}: a arquitetura dele fixa o banco em 8 filtros por camada, os dois braços acrescentam exatamente 24, e o deslocamento que a intervenção evita não pode ocorrer. Semente cuja base degenerou (F$\beta$=0) sai, porque $\Delta$ a partir de zero não pode ser negativo. \textbf{O desfecho primário vem com três testes}: t pareado, teste de sinais e permutação exata de sinais. Com 3 e 5 sementes a normalidade das diferenças não é verificável, e o t sozinho não autorizaria conclusão; onde os três concordam, como no conjunto de parasitas, a conclusão não depende da suposição. \textbf{ATENÇÃO ao pior caso da conjuntivite}: o valor de 0,0480 é bruto e \textbf{não sobrevive} à análise de sensibilidade que casa o número de candidatos entre os braços, na qual passa a 0,1032. O gerador de contorno entregou UM único candidato naquela semente, e mínimo de um sorteio contra mínimo de oito não é comparação. Esse desfecho \textbf{não deve ser lido como significativo}.
\normalsize


No conjunto de parasitas o efeito é claro: ganho médio de 0,0772 de Fβ, com intervalo de confiança entre 0,0393 e 0,1151, valor p de 0,0013, e vantagem em todas as dez repetições. O resultado sobrevive à correção de Benjamini-Hochberg aplicada sobre os três conjuntos de dados. Os dois desfechos secundários, declarados no pré-registro com o mecanismo escrito antes da execução, acompanham: o pior caso melhora 0,2153 e a dispersão entre candidatos cai 0,1091.

O comportamento nos outros dois conjuntos merece atenção. Na conjuntivite o efeito vai na mesma direção, com vantagem em quatro das cinco repetições, mas sem significância no desfecho primário. No BraTS o efeito desaparece, e esse resultado nulo é previsto pelo mecanismo: a arquitetura empregada nesse conjunto fixa o banco em oito filtros por camada, de modo que ambos os braços acrescentam exatamente o mesmo número de filtros e o deslocamento que a intervenção evita não pode ocorrer. O conjunto funciona, assim, como controle do mecanismo.

Duas limitações do gerador foram medidas e devem acompanhar qualquer leitura do resultado. A primeira é que ele exige uma predição não trivial: quando a máscara prevista é vazia ou cobre a imagem inteira, não há contorno e o conjunto de candidatos fica vazio, enquanto o gerador por superpixels funciona em qualquer situação. A segunda é que os candidatos gerados apresentam fração de objeto próxima de um, isto é, caem no interior do objeto verdadeiro e não sobre sua fronteira, porque o modelo subsegmenta e o contorno previsto se situa dentro do contorno real. O ganho ocorre apesar disso, o que sugere que parte do efeito decorre de o candidato ser pequeno e localizado, acrescentando menos filtros ao banco, e não apenas de sua posição. Separar essas duas contribuições exige um experimento adicional, discutido no Capítulo \ref{cap_5}.

\section{Resultados de natureza metodológica}

Três resultados não dizem respeito ao desempenho dos métodos, mas à validade das medições, e por isso são relatados separadamente.

\subsection{Reprodutibilidade do treinamento}

O treinamento do codificador não era reprodutível entre execuções. Duas chamadas idênticas, com os mesmos arquivos de marcadores e no mesmo processo, produziam pesos diferentes quando o número de candidatos a filtro ultrapassava o limite da arquitetura: a diferença máxima observada nos pesos da segunda camada foi de 0,134. Abaixo desse limite o procedimento devolve os próprios candidatos sem agrupá-los, e por isso é estável; acima dele o agrupamento precisa escolher entre muitos candidatos, e a ordem de redução das somas em bibliotecas de álgebra linear com múltiplas linhas de execução altera o resultado no último dígito, o que basta para o agrupamento convergir para outra configuração. O efeito se amplifica entre camadas, já que cada camada recebe como entrada as ativações da anterior.

Fixando as bibliotecas em uma única linha de execução, a diferença passa a ser exatamente zero nos dois regimes. Todas as campanhas posteriores a essa descoberta foram executadas sob essa restrição.

\subsection{Defeitos que invalidaram medições anteriores}

Três defeitos de implementação foram encontrados e corrigidos, e os números produzidos antes das correções não são válidos. O procedimento de extração de atributos para os critérios de diversidade assumia imagens de três bandas, o que impedia sua execução no conjunto BraTS, de banda única. A leitura de imagens assumia extensão fixa, o que impedia a execução no conjunto de conjuntivite. E o filtro de área do pós-processamento estava configurado com a faixa adequada ao conjunto de parasitas, cujos objetos são pequenos, o que eliminava toda componente prevista nas imagens de conjuntivite: nenhum dos objetos examinados cabia na faixa, e o Fβ mediano passou de 0,024 para 0,594 após a correção.

Esse último caso tem interesse além do defeito pontual. Ele indica que os parâmetros de pós-processamento publicados não transferem entre domínios, e que aplicá-los sem reavaliação pode produzir resultados sem relação com o método avaliado.

\subsection{Falsos positivos capturados pelo protocolo}

O protocolo de pré-registro, descrito no Capítulo \ref{cap_3}, capturou quatro resultados que não se sustentaram. Dois deles haviam atingido significância em análises exploratórias e foram submetidos a confirmatórios com repetições inéditas: um apresentava valor p de 0,0016 e passou a 0,586, outro apresentava 0,082 e passou a 0,567. Um terceiro desapareceu quando o número de repetições foi completado, com o valor p passando de 0,0265 para 0,1385. O quarto foi falseado não apenas na hipótese como no mecanismo proposto, pelo critério declarado antes da execução.

\begin{table}[htb]
\centering
\caption{Composição do registro de execuções.}
\label{tab_proveniencia}
% Proveniência disponível por família de experimento
% GERADO por flim_al/tabelas.py em 2026-10-08T14:01:07+00:00
% NAO EDITE A MAO: a proxima geracao sobrescreve, e o numero editado deixa de bater com evidencia/execucoes/.
% Proveniencia: proveniencia_registro.proveniencia.json
\begin{tabular}{lrrr}
\toprule
família & execuções & sem seed & sem commit \\
\midrule
ganho\_marginal & 2448 & 0\% & 0\% \\
al\_corrected\_results & 2316 & 42\% & 100\% \\
tabela\_k\_por\_modelo & 1840 & 0\% & 0\% \\
final\_comparison\_v3 & 1029 & 4\% & 100\% \\
benchmark\_conjunctiva & 588 & 68\% & 100\% \\
sel\_cs\_grande & 540 & 0\% & 100\% \\
al\_backprop\_results & 492 & 20\% & 100\% \\
paper\_selection\_pac2 & 432 & 0\% & 100\% \\
paper\_selection\_POOL\_COM\_VAZAMENTO & 360 & 0\% & 100\% \\
paper\_selection\_v2 & 324 & 0\% & 100\% \\
paper\_selection\_v3 & 324 & 0\% & 100\% \\
paper\_selection\_rand & 321 & 0\% & 100\% \\
densidade & 262 & 0\% & 0\% \\
cl\_plasticidade & 254 & 0\% & 0\% \\
sel\_degen & 216 & 0\% & 100\% \\
sel\_medoide & 216 & 0\% & 100\% \\
sel\_proposto & 216 & 0\% & 100\% \\
final\_comparison\_v2 & 210 & 10\% & 100\% \\
il\_noc & 207 & 0\% & 0\% \\
clique & 183 & 0\% & 0\% \\
paper\_selection\_coreset & 162 & 0\% & 100\% \\
onde\_marcar & 150 & 0\% & 0\% \\
final\_comparison & 147 & 57\% & 100\% \\
al\_encoder\_results\_acq\_comparison & 144 & 44\% & 100\% \\
orcamento & 144 & 100\% & 100\% \\
artigo\_vs\_regiao & 130 & 0\% & 0\% \\
al\_curve & 84 & 100\% & 100\% \\
brats & 84 & 25\% & 100\% \\
al\_encoder\_results & 70 & 70\% & 100\% \\
al\_encoder\_results\_dt & 62 & 47\% & 100\% \\
paper\_selection\_brats & 54 & 0\% & 100\% \\
paper\_selection\_uB\_POOL\_COM\_VAZAMENTO & 54 & 0\% & 100\% \\
al\_bald\_results & 48 & 44\% & 100\% \\
al\_flim\_curve & 42 & 100\% & 100\% \\
curva & 24 & 100\% & 100\% \\
al\_region\_results & 16 & 44\% & 100\% \\
\bottomrule
\end{tabular}

\end{table}

% Nota de proveniencia_registro. GERADA por flim_al/tabelas.py.
\footnotesize Diagnóstico, não resultado. Linha com 100\% sem seed não sustenta afirmação sobre variância entre execuções.
\normalsize


\section{Síntese}

Os resultados sustentam a seguinte leitura. Critérios de \ac{AL} não apresentam superioridade sobre a escolha aleatória no \ac{FLIM}, seja na seleção de imagens, seja na seleção de regiões, nas condições avaliadas. Existe margem considerável na escolha da região, da ordem de 0,13 a 0,30 de Fβ, e ela não é alcançada pelos critérios conhecidos. A razão é estrutural e tem duas partes: a anotação adicional reestima o extrator de atributos em vez de refiná-lo, o que torna o acréscimo uma intervenção de risco; e o procedimento que gera os candidatos quase não oferece a anotação que os marcadores dos especialistas mostram ter valor.

A intervenção sobre o gerador, derivada desse diagnóstico, produz ganho mensurável em um dos três conjuntos de dados, com o comportamento nos demais explicado pelo mesmo mecanismo. O Capítulo \ref{cap_5} discute o que esses resultados autorizam afirmar e o que permanece em aberto.
